\chapter{Electrodesiccation and Curettage}

\section*{Overview}
Electrodesiccation and curettage (ED\&C) is a dermatologic procedure in which a skin lesion is scraped away with a curette and the base is treated with a low-voltage electric current to destroy remaining tissue and control bleeding. It is used mainly to treat basal cell carcinoma and superficial skin lesions.
\section*{Alternative Names}
ED\&C; curettage and electrodesiccation; curettage and cautery; scrape and burn
\section*{Principle}
The curette removes the soft, friable tumor tissue, which differs in consistency from normal skin, and electrodesiccation destroys residual malignant cells along the margins and base. Repeated cycles extend treatment depth and help achieve tumor clearance.
\section*{History}
Curettage with electrosurgical destruction was developed in the early 20th century as a simple method for treating skin cancers and remains widely used for low-risk basal cell carcinoma.
\section*{Why It Is Done}
Ordered for small, low-risk basal cell carcinomas, superficial squamous cell carcinoma in situ, actinic keratoses, seborrheic keratoses, and other benign or premalignant skin lesions when excision or Mohs surgery is not required.
\section*{Is This a Primary or Secondary Test or Procedure}
Primary therapeutic procedure for selected low-risk skin lesions.
\section*{How It Is Performed}
The area is cleaned and anesthetized with local anesthetic. The lesion is scraped with a dermal curette, the base is treated with an electrosurgical unit to desiccate the tissue, and the cycle is repeated two to three times. The wound is dressed and heals by secondary intention.
\section*{Preparation}
Informed consent and review of the patient's bleeding risk and any implanted electronic devices, since electrosurgery may interfere with pacemakers or other implants.
\section*{Contraindications and Precautions}
Precautions include lesions near the eye or nose (risk of poor cosmetic and functional results), large or deeply invasive tumors, and conditions with high recurrence risk where Mohs surgery or excision is preferred.
\section*{Risks}
Pain, bleeding, infection, scarring, hypopigmentation, and recurrence of the treated lesion, particularly with larger or more aggressive tumors.
\section*{Aftercare and Recovery}
The wound is kept clean and dressed until healing occurs over several weeks. The patient is advised on signs of infection, and the treated area heals with a flat scar.
\section*{Outcome and Follow-up}
Successful treatment is indicated by complete removal of the lesion with healing and no local recurrence at follow-up. Recurrence may require a different treatment modality.
\section*{Expected Results}
Complete elimination of the lesion with a cosmetically acceptable scar and no clinical evidence of recurrence.
\section*{Follow-up}
Patients are reassessed for healing and recurrence, and periodic full skin examinations are recommended for those with a history of skin cancer. Mohs surgery or excision is considered for recurrent lesions.
\section*{References}
\begin{enumerate}
  \item MedlinePlus. Skin cancer - treatments. U.S. National Library of Medicine; 2025.
  \item Current Medical Diagnosis and Treatment. McGraw-Hill; 2025.
\end{enumerate}

\clearpage
